\chapter{Catalogs and Commits}
\label{ch:catalogs}

\begin{chapterintro}
The catalog is the one component that turns a table name into the location of its
current metadata --- and the point at which a commit becomes atomic. Writing new
metadata is not a commit until the catalog's pointer is swapped to it. When two
writers race, one wins the swap and the other retries. This chapter covers what
the Hive Metastore and REST catalogs do, how the atomic swap works, and how to
tune retries so a busy table does not fail commits under contention.
\end{chapterintro}

\section{What a catalog does}
\tbd

\section{The Hive Metastore catalog}
\tbd

\section{The REST catalog}
\tbd

\section{The atomic pointer swap}
\tbd

\section{Optimistic concurrency and retries}
\tbd

\section{Catalog configuration in Spark}
\tbd

\section{Summary}
\tbd

\exercisehead
\begin{exercises}
\item Run two writers into the same table. Observe the retry, then find the
      winning and losing commit attempts in the snapshot log.
\end{exercises}
